\chapter{Lenzen, beelden en het oog}\label{ch:g11:lenses-and-eye}

Elk scherp beeld --- op een telefoonsensor, een bioscoopdoek of achter in
je eigen oog --- ontstaat door dezelfde truc: een gekromd doorzichtig
midden breekt stralen (\cref{ch:g10:refraction}) zodat alle stralen die
uit één punt vertrekken elkaar in een ander punt weer ontmoeten. Dit
hoofdstuk temt die truc: één \omterm{def:g11:lenses-and-eye:lens}{lens}, drie stralen, één formule --- en richt
hem daarna op fototoestellen, brillen en de levende \omterm{def:g11:lenses-and-eye:lens}{lens} die deze zin
leest.

\section{Convergerende en divergerende lenzen}

\begin{definition}[Dunne lens]\label{def:g11:lenses-and-eye:lens}
Een \emph{lens}\index{lens} is een doorzichtig midden dat door twee
oppervlakken wordt begrensd, waarvan er minstens één gekromd is. Is de
lens in het midden dikker dan aan de rand, dan is ze
\emph{convergerend}\index{convergerende lens}: evenwijdige stralen
worden naar elkaar toe gebogen; is ze aan de rand dikker, dan is ze
\emph{divergerend}\index{divergerende lens} en lopen ze uiteen. Een
\emph{dunne}\index{dunne lens} lens heeft een verwaarloosbare dikte: we
tekenen haar als een lijnstuk met pijlpunten naar buiten (convergerend)
of naar binnen (divergerend). Haar middelpunt $O$ is het \emph{optische
middelpunt}\index{optisch middelpunt}; de lijn door $O$ loodrecht op de
lens is de \emph{hoofdas}\index{hoofdas}.
\end{definition}

\begin{definition}[Brandpunten en brandpuntsafstand]\label{def:g11:lenses-and-eye:focal}
Stralen die evenwijdig aan de \omterm{def:g11:lenses-and-eye:lens}{hoofdas} op een \omterm{def:g11:lenses-and-eye:lens}{convergerende lens} vallen,
snijden de as alle in één punt erachter: het \emph{beeldbrandpunt}%
\index{brandpunt} $F'$; omgekeerd verlaten stralen die uit het
\emph{voorwerpsbrandpunt} $F$ vertrekken (ervoor, met $OF = OF'$) de \omterm{def:g11:lenses-and-eye:lens}{lens}
evenwijdig aan de as. De \emph{brandpuntsafstand}\index{brandpuntsafstand}
$f' = \overline{OF'}$ is positief voor een \omterm{def:g11:lenses-and-eye:lens}{convergerende lens}; door een
\omterm{def:g11:lenses-and-eye:lens}{divergerende lens} verlaten evenwijdige stralen de \omterm{def:g11:lenses-and-eye:lens}{lens} alsof ze uit een
punt $F'$ \emph{vóór} de \omterm{def:g11:lenses-and-eye:lens}{lens} komen: $f' < 0$.
\end{definition}

\begin{definition}[Lenssterkte]\label{def:g11:lenses-and-eye:vergence}
De \emph{lenssterkte}\index{lenssterkte} van een \omterm{def:g11:lenses-and-eye:lens}{lens} met
\omterm{def:g11:lenses-and-eye:focal}{brandpuntsafstand} $f'$ (in \omterm{def:g10:orders-of-magnitude:unit}{meter}) is $C = 1/f'$, gemeten in
\emph{dioptrie}\index{dioptrie} (symbool D): $1~\text{D} =
\qty{1}{m^{-1}}$; divergerende lenzen hebben $C < 0$. Opticiens
schrijven in dioptrie voor: de sterkten van \omterm{def:g11:lenses-and-eye:lens}{dunne} lenzen die elkaar raken
tellen gewoon op (hier aangenomen; afgeleid in het volume van
bachelorjaar 1).
\end{definition}

\begin{example}[Een recept lezen]\label{ex:g11:lenses-and-eye:vergence}
Een \omterm{def:g11:lenses-and-eye:lens}{lens} met $f' = \qty{0.50}{m}$ heeft $C = 1/0.50 = +2.0$~D:
\omterm{def:g11:lenses-and-eye:lens}{convergerend}. Een voorschrift van $-4.0$~D betekent
$f' = \qty{-0.25}{m}$: \omterm{def:g11:lenses-and-eye:lens}{divergerend}, met \omterm{def:g11:lenses-and-eye:focal}{brandpuntsafstand} \qty{25}{cm}.
\end{example}

\section{Het beeld construeren}

\begin{method}[De drie constructiestralen]\label{met:g11:lenses-and-eye:rays}
Zij $AB$ een voorwerp dat in $A$ loodrecht op de as staat. Van alle
stralen die uit de top $B$ vertrekken, zijn er drie waarvan de weg bekend
is:
\begin{enumerate}
  \item de straal door het \omterm{def:g11:lenses-and-eye:lens}{optische middelpunt} $O$ gaat rechtdoor;
  \item de straal evenwijdig aan de as gaat na de \omterm{def:g11:lenses-and-eye:lens}{lens} door $F'$;
  \item de straal door $F$ verlaat de \omterm{def:g11:lenses-and-eye:lens}{lens} evenwijdig aan de as.
\end{enumerate}
Twee ervan volstaan: hun snijpunt is het beeldpunt $B'$, en het beeld
$A'B'$ staat in $A'$ loodrecht op de as. Lopen de uittredende stralen
uiteen, dan snijden hun \emph{achterwaartse} verlengden (gestippeld)
elkaar in een \omterm{def:g11:lenses-and-eye:images}{virtueel} beeldpunt.
\end{method}

\begin{omfigure}
\begin{tikzpicture}[scale=0.55]
  \draw[-Stealth] (-6.4,0) -- (6.4,0);
  \draw[Stealth-Stealth, omDef, very thick] (0,-2.4) -- (0,2.4);
  \fill (0,0) circle (1.5pt) node[below right, font=\small] {$O$};
  \fill (-2,0) circle (1.5pt) node[below left, font=\small] {$F$};
  \fill (2,0) circle (1.5pt) node[below right, font=\small] {$F'$};
  \draw[omMeth] (-4,1.2) -- (0,1.2) -- (4,-1.2);
  \draw[omMeth] (-4,1.2) -- (4,-1.2);
  \draw[omMeth] (-4,1.2) -- (0,-1.2) -- (4,-1.2);
  \draw[-Stealth, very thick] (-4,0) -- (-4,1.2);
  \node[above, font=\small] at (-4,1.2) {$B$};
  \node[below, font=\small] at (-4,0) {$A$};
  \draw[-Stealth, omThm, very thick] (4,0) -- (4,-1.2);
  \node[below, font=\small] at (4,-1.2) {$B'$};
  \node[above, font=\small] at (4,0) {$A'$};
\end{tikzpicture}

{\small De drie constructiestralen: voorwerp $AB$ voorbij $F$, beeld
$A'B'$ \omterm{def:g11:lenses-and-eye:images}{reëel} en omgekeerd (hier is $\overline{OA} = -2f'$: even groot,
$\gamma = -1$).}
\end{omfigure}

\begin{definition}[Reële en virtuele beelden]\label{def:g11:lenses-and-eye:images}
Een beeld is \emph{reëel}\index{reëel beeld} wanneer de uittredende
stralen er werkelijk doorheen gaan: een scherm dat je daar plaatst vangt
het op; het is \emph{virtueel}\index{virtueel beeld} wanneer alleen hun
achterwaartse verlengden elkaar snijden: geen scherm vangt het op, maar
een oog dat in de \omterm{def:g11:lenses-and-eye:lens}{lens} kijkt ziet het uitstekend. Een \omterm{def:g11:lenses-and-eye:lens}{convergerende lens}
geeft van elk voorwerp voorbij $F$ een reëel, omgekeerd beeld; schuif het
voorwerp binnen de \omterm{def:g11:lenses-and-eye:focal}{brandpuntsafstand} en het beeld wordt virtueel,
rechtop en vergroot.
\end{definition}

\begin{omfigure}
\begin{tikzpicture}[scale=0.5]
  \draw[-Stealth] (-7.2,0) -- (4.6,0);
  \draw[Stealth-Stealth, omDef, very thick] (0,-2.2) -- (0,3.6);
  \fill (0,0) circle (1.5pt) node[below right, font=\small] {$O$};
  \fill (-2,0) circle (1.5pt) node[below, font=\small] {$F$};
  \fill (2,0) circle (1.5pt) node[below, font=\small] {$F'$};
  \draw[omMeth] (-1.5,0.8) -- (3,-1.6);
  \draw[omMeth] (-1.5,0.8) -- (0,0.8) -- (4,-0.8);
  \draw[omMeth, dashed] (0,0.8) -- (-6,3.2);
  \draw[omMeth, dashed] (-1.5,0.8) -- (-6,3.2);
  \draw[-Stealth, very thick] (-1.5,0) -- (-1.5,0.8);
  \node[below, font=\small] at (-1.5,0) {$A$};
  \node[above right, font=\small] at (-1.4,0.75) {$B$};
  \draw[-Stealth, omThm, very thick, dashed] (-6,0) -- (-6,3.2);
  \node[below, font=\small] at (-6,0) {$A'$};
  \node[above, font=\small] at (-6,3.2) {$B'$};
\end{tikzpicture}

{\small Voorwerp binnen de \omterm{def:g11:lenses-and-eye:focal}{brandpuntsafstand} (de loepstand): de
achterwaartse verlengden bouwen een \omterm{def:g11:lenses-and-eye:images}{virtueel}, rechtopstaand, vergroot
beeld $A'B'$ op.}
\end{omfigure}

\section{Eén formule: de lenzenformule}

Constructies laten zien \emph{waar} het beeld ligt; een formule berekent
het. Kies de as in de richting van het licht en gebruik
\emph{algebraïsche maten}: $\overline{OA}$ is de coördinaat van $A$ ten
opzichte van $O$ (negatief voor een \omterm{def:g11:lenses-and-eye:images}{reëel} voorwerp ervoor); hoogten
$\overline{AB}$ tellen positief naar boven.

\begin{theorem}[Lenzenformule voor dunne lenzen]\label{thm:g11:lenses-and-eye:conjugate}
Door een \omterm{def:g11:lenses-and-eye:lens}{dunne lens} met \omterm{def:g11:lenses-and-eye:focal}{brandpuntsafstand} $f'$ wordt een voorwerpspunt
$A$ op de as afgebeeld op het punt $A'$ met
\[
\frac{1}{\overline{OA'}} - \frac{1}{\overline{OA}} = \frac{1}{f'} = C .
\]
\end{theorem}
\admitted

\begin{remark}
De betrekking volgt uit de wet van Snellius, toegepast op stralen dicht
bij de as; de afleiding --- en de grenzen van het model van de \omterm{def:g11:lenses-and-eye:lens}{dunne lens}
--- worden eerlijk gemaakt in het volume van bachelorjaar 1.
\end{remark}

\begin{definition}[Vergroting]\label{def:g11:lenses-and-eye:magnification}
De \emph{vergroting}\index{vergroting} van het beeld is
\[
\gamma = \frac{\overline{A'B'}}{\overline{AB}}
       = \frac{\overline{OA'}}{\overline{OA}} .
\]
$\abs{\gamma}$ vergelijkt de afmetingen, het teken de standen
($\gamma < 0$: omgekeerd); het teken van $\overline{OA'}$ verraadt de
aard (positief: \omterm{def:g11:lenses-and-eye:images}{reëel}, achter de \omterm{def:g11:lenses-and-eye:lens}{lens}; negatief: \omterm{def:g11:lenses-and-eye:images}{virtueel}, ervoor).
\end{definition}

\begin{example}[De formule aan het werk]\label{ex:g11:lenses-and-eye:conjugate}
Een voorwerp staat \qty{30}{cm} voor een \omterm{def:g11:lenses-and-eye:lens}{convergerende lens} met
$f' = \qty{10}{cm}$: $\overline{OA} = \qty{-30}{cm}$, dus
\[
\frac{1}{\overline{OA'}} = \frac{1}{10} - \frac{1}{30} = \frac{1}{15},
\qquad \overline{OA'} = +\qty{15}{cm},
\qquad \gamma = \frac{15}{-30} = -0.50 :
\]
een \omterm{def:g11:lenses-and-eye:images}{reëel beeld} \qty{15}{cm} achter de \omterm{def:g11:lenses-and-eye:lens}{lens}, omgekeerd en half zo groot
--- precies wat de drie stralen tekenen.
\end{example}

\section{Het oog}

\begin{definition}[Het oog als optisch systeem]\label{def:g11:lenses-and-eye:eye}
Optisch gezien is het oog een \omterm{def:g11:lenses-and-eye:lens}{convergerende lens} tegenover een scherm:
het \emph{hoornvlies}\index{hoornvlies} en de \emph{ooglens}\index{ooglens}
werken samen als één \omterm{def:g11:lenses-and-eye:lens}{convergerende lens} met instelbare
\omterm{def:g11:lenses-and-eye:focal}{brandpuntsafstand}; het \emph{netvlies}\index{netvlies}, ongeveer
\qty{17}{mm} erachter, is het vaste scherm. Scherpstellen gebeurt door de
\omterm{def:g11:lenses-and-eye:lens}{lens} samen te knijpen, niet door haar te verschuiven zoals in een
fototoestel: \emph{accommodatie}\index{accommodatie} is de verandering
van $f'$ die de kringspieren teweegbrengen. Het dichtstbijzijnde punt dat
bij volledige accommodatie scherp wordt gezien is het
\emph{nabijpunt}\index{nabijpunt} --- ongeveer \qty{25}{cm} voor het
\emph{normale oog}\index{normaal oog} --- en het verste, in rust, het
\emph{vertepunt}\index{vertepunt} (oneindig voor het normale oog).
\end{definition}

\begin{omfigure}
\begin{tikzpicture}[scale=0.85]
  \draw[thick] (0,0) circle (1.1);
  \draw[omDef, very thick] (-0.95,-0.5) -- (-0.95,0.5);
  \draw (-2.6,0.35) -- (-0.95,0.35) -- (1.1,0);
  \draw (-2.6,0) -- (1.1,0);
  \draw (-2.6,-0.35) -- (-0.95,-0.35) -- (1.1,0);
  \fill (1.1,0) circle (1.5pt) node[right, font=\small] {netvlies};
  \node[font=\small] at (-0.7,-1.5) {in rust: ver voorwerp};
\end{tikzpicture}\qquad
\begin{tikzpicture}[scale=0.85]
  \draw[thick] (0,0) circle (1.1);
  \draw[omDef, line width=2.6pt] (-0.95,-0.5) -- (-0.95,0.5);
  \fill (-2.6,0) circle (1.5pt) node[above, font=\small] {$B$};
  \draw (-2.6,0) -- (-0.95,0.4) -- (1.1,0);
  \draw (-2.6,0) -- (1.1,0);
  \draw (-2.6,0) -- (-0.95,-0.4) -- (1.1,0);
  \node[font=\small] at (-0.7,-1.5) {geaccommodeerd: nabij voorwerp};
\end{tikzpicture}

{\small \omterm{def:g11:lenses-and-eye:eye}{Accommodatie}: om een nabij voorwerp op hetzelfde \omterm{def:g11:lenses-and-eye:eye}{netvlies} scherp
te krijgen, bolt de \omterm{def:g11:lenses-and-eye:eye}{ooglens} op --- haar \omterm{def:g11:lenses-and-eye:vergence}{lenssterkte} neemt toe.}
\end{omfigure}

\begin{example}[De sterkte van je oog]\label{ex:g11:lenses-and-eye:eyerange}
Het \omterm{def:g11:lenses-and-eye:eye}{netvlies} zit op $\overline{OA'} = +\qty{0.017}{m}$. Ver voorwerp
($1/\overline{OA} \approx 0$): $C = 1/0.017 \approx 59$~D. \omterm{def:g11:lenses-and-eye:eye}{Nabijpunt}
($\overline{OA} = \qty{-0.25}{m}$): $C = 1/0.017 + 1/0.25 \approx 63$~D.
Lezen kost ongeveer $4$~D: het oog is een \omterm{def:g11:lenses-and-eye:lens}{lens} van $59$~D met een
ingebouwde fijnregeling van $+4$~D.
\end{example}

\section{Afwijkingen, correcties en de loep}

\begin{definition}[Bijziendheid en verziendheid]\label{def:g11:lenses-and-eye:defects}
Een \emph{bijziend}\index{bijziendheid} oog is te sterk \omterm{def:g11:lenses-and-eye:lens}{convergerend} voor
zijn lengte: evenwijdige stralen komen \emph{vóór} het \omterm{def:g11:lenses-and-eye:eye}{netvlies} samen,
verre voorwerpen worden wazig en het \omterm{def:g11:lenses-and-eye:eye}{vertepunt} ligt op eindige afstand;
een \omterm{def:g11:lenses-and-eye:lens}{divergerend} brillenglas verhelpt dat. Een
\emph{verziend}\index{verziendheid} oog convergeert niet genoeg: nabije
voorwerpen zouden achter het \omterm{def:g11:lenses-and-eye:eye}{netvlies} scherp komen en het \omterm{def:g11:lenses-and-eye:eye}{nabijpunt}
schuift weg; een \omterm{def:g11:lenses-and-eye:lens}{convergerende lens} verhelpt dat. Met de jaren verstijft
de \omterm{def:g11:lenses-and-eye:eye}{ooglens} (\emph{ouderdomsverziendheid}\index{ouderdomsverziendheid}) en
krimpt de \omterm{def:g11:lenses-and-eye:eye}{accommodatie} --- met hetzelfde geneesmiddel: een leesbril.
\end{definition}

\begin{example}[Een voorschrift voor een bijziend oog]\label{ex:g11:lenses-and-eye:correction}
Een \omterm{def:g11:lenses-and-eye:defects}{bijziend} oog heeft zijn \omterm{def:g11:lenses-and-eye:eye}{vertepunt} op \qty{50}{cm}: het corrigerende
glas moet verre voorwerpen \emph{op het \omterm{def:g11:lenses-and-eye:eye}{vertepunt}} laten zien ---
voorwerp op oneindig, beeld op $\overline{OA'} = \qty{-0.50}{m}$
(\omterm{def:g11:lenses-and-eye:images}{virtueel}, ervoor). Dan geldt $1/f' = 1/\overline{OA'}$:
$f' = \qty{-0.50}{m}$, $C = -2.0$~D; het ontspannen oog kijkt naar dat
virtuele beeld en ziet het scherp.
\end{example}

\begin{omfigure}
\begin{tikzpicture}[scale=0.85]
  \draw[thick] (0,0) circle (1.1);
  \draw[omDef, very thick] (-0.95,-0.5) -- (-0.95,0.5);
  \draw (-2.7,0.35) -- (-0.95,0.35) -- (0.55,0) -- (1.1,-0.13);
  \draw (-2.7,0) -- (1.1,0);
  \draw (-2.7,-0.35) -- (-0.95,-0.35) -- (0.55,0) -- (1.1,0.13);
  \fill[omThm] (0.55,0) circle (1.5pt);
  \node[font=\small] at (-0.7,-1.5) {bijziend: scherp vóór het netvlies};
\end{tikzpicture}\qquad
\begin{tikzpicture}[scale=0.85]
  \draw[thick] (0,0) circle (1.1);
  \draw[omDef, very thick] (-0.95,-0.5) -- (-0.95,0.5);
  \draw[omMeth, very thick, {Stealth[reversed]}-{Stealth[reversed]}]
    (-1.9,-0.55) -- (-1.9,0.55);
  \draw (-2.8,0.3) -- (-1.9,0.3) -- (-0.95,0.42) -- (1.1,0);
  \draw (-2.8,0) -- (1.1,0);
  \draw (-2.8,-0.3) -- (-1.9,-0.3) -- (-0.95,-0.42) -- (1.1,0);
  \node[font=\small] at (-0.7,-1.5) {gecorrigeerd: divergerende lens};
\end{tikzpicture}

{\small Een \omterm{def:g11:lenses-and-eye:defects}{bijziend} oog stelt evenwijdige stralen vóór het \omterm{def:g11:lenses-and-eye:eye}{netvlies}
scherp (rode punt); een \omterm{def:g11:lenses-and-eye:lens}{divergerende lens} legt het \omterm{def:g11:lenses-and-eye:focal}{brandpunt} op het
\omterm{def:g11:lenses-and-eye:eye}{netvlies}.}
\end{omfigure}

\begin{definition}[De loep]\label{def:g11:lenses-and-eye:magnifier}
Een \emph{loep}\index{loep} is een \omterm{def:g11:lenses-and-eye:lens}{convergerende lens} met een korte
\omterm{def:g11:lenses-and-eye:focal}{brandpuntsafstand}, gebruikt met het voorwerp \emph{binnen} die
\omterm{def:g11:lenses-and-eye:focal}{brandpuntsafstand}: het beeld is \omterm{def:g11:lenses-and-eye:images}{virtueel}, rechtop en vergroot
(\cref{def:g11:lenses-and-eye:images}), en het oog bekijkt het op een
comfortabele afstand.
\end{definition}

\begin{example}[Een juweliersloep]\label{ex:g11:lenses-and-eye:magnifier}
\omterm{def:g11:lenses-and-eye:magnifier}{Loep} met $f' = \qty{5.0}{cm}$, steen op \qty{4.0}{cm}:
$1/\overline{OA'} = 1/5.0 - 1/4.0 = -1/20$, dus
$\overline{OA'} = \qty{-20}{cm}$ en $\gamma = (-20)/(-4.0) = +5.0$: een
\omterm{def:g11:lenses-and-eye:images}{virtueel}, rechtopstaand beeld, vijf keer zo groot, op een comfortabele
\qty{20}{cm} van de \omterm{def:g11:lenses-and-eye:lens}{lens}.
\end{example}

\section{Oefeningen}

\begin{exercise}[$\star$]\label{exo:g11:lenses-and-eye:1}
Bereken de \omterm{def:g11:lenses-and-eye:vergence}{lenssterkten} voor $f' = \qty{50}{cm}$, $\qty{-20}{cm}$ en
$\qty{4.0}{mm}$, en daarna de \omterm{def:g11:lenses-and-eye:focal}{brandpuntsafstand} van een \omterm{def:g11:lenses-and-eye:lens}{lens} van
$+8.0$~D.
\end{exercise}

\begin{exercise}[$\star$]\label{exo:g11:lenses-and-eye:2}
Een \omterm{def:g11:lenses-and-eye:lens}{lens} is aan de randen dikker dan in het midden. \omterm{def:g11:lenses-and-eye:lens}{Convergerend} of
\omterm{def:g11:lenses-and-eye:lens}{divergerend}? Kan ze als \omterm{def:g11:lenses-and-eye:magnifier}{loep} dienen? Wat zie je als je haar boven een
bedrukte bladzijde houdt?
\end{exercise}

\begin{exercise}[$\star$]\label{exo:g11:lenses-and-eye:3}
Een \omterm{def:g11:lenses-and-eye:lens}{convergerende lens} heeft $f' = \qty{10}{cm}$; een voorwerp staat
\qty{20}{cm} ervoor. Construeer het beeld met twee van de drie stralen en
toets daarna plaats en grootte met de lenzenformule.
\end{exercise}

\begin{exercise}[$\star$]\label{exo:g11:lenses-and-eye:4}
Een voorwerp staat \qty{24}{cm} voor een \omterm{def:g11:lenses-and-eye:lens}{convergerende lens} met
$f' = \qty{8.0}{cm}$. Bereken $\overline{OA'}$ en $\gamma$, en beschrijf
het beeld (aard, stand, grootte).
\end{exercise}

\begin{exercise}[$\star$]\label{exo:g11:lenses-and-eye:5}
In het oog: wat vormt de \omterm{def:g11:lenses-and-eye:lens}{convergerende lens}? Wat speelt het scherm? Wat
verandert er bij \omterm{def:g11:lenses-and-eye:eye}{accommodatie} --- en wat kan er niet veranderen?
Definieer het \omterm{def:g11:lenses-and-eye:eye}{nabijpunt} van het \omterm{def:g11:lenses-and-eye:eye}{normale oog}.
\end{exercise}

\begin{exercise}[$\star\star$]\label{exo:g11:lenses-and-eye:6}
Een objectief heeft $f' = \qty{50}{mm}$. Waar zit de sensor bij
scherpstellen op een ver landschap? Bereken de afstand lens--sensor voor
een onderwerp op \qty{3.0}{m}, en de verplaatsing tussen beide
instellingen.
\end{exercise}

\begin{exercise}[$\star\star$]\label{exo:g11:lenses-and-eye:7}
Een \omterm{def:g11:lenses-and-eye:magnifier}{loep} heeft $f' = \qty{6.0}{cm}$; een postzegel ligt \qty{5.0}{cm}
eronder. Plaats, aard, stand en \omterm{def:g11:lenses-and-eye:magnification}{vergroting} van het beeld?
\end{exercise}

\begin{exercise}[$\star\star$]\label{exo:g11:lenses-and-eye:8}
Een voorwerp en een scherm staan tegenover elkaar; een \omterm{def:g11:lenses-and-eye:lens}{convergerende lens}
ertussen geeft een scherp beeld wanneer beide \qty{40}{cm} van de \omterm{def:g11:lenses-and-eye:lens}{lens}
staan. Bereken $f'$ en $\gamma$. Waarom moet dit symmetrische beeld
precies even groot zijn als het voorwerp?
\end{exercise}

\begin{exercise}[$\star\star$]\label{exo:g11:lenses-and-eye:9}
Een voorwerp staat \qty{50}{cm} voor een \omterm{def:g11:lenses-and-eye:lens}{divergerende lens} met
$f' = \qty{-25}{cm}$. Bereken $\overline{OA'}$ en $\gamma$; beschrijf het
beeld. Waarom lijken de ogen van een bijziende iets kleiner door zijn
bril?
\end{exercise}

\begin{exercise}[$\star\star$]\label{exo:g11:lenses-and-eye:10}
Een \omterm{def:g11:lenses-and-eye:defects}{bijziend} oog heeft zijn \omterm{def:g11:lenses-and-eye:eye}{vertepunt} op \qty{40}{cm}. Welke \omterm{def:g11:lenses-and-eye:vergence}{lenssterkte}
moet het corrigerende glas hebben, en waar legt het het beeld van een
verre berg? Waarom wordt die berg dan moeiteloos gezien?
\end{exercise}

\begin{exercise}[$\star\star$]\label{exo:g11:lenses-and-eye:11}
Een lezer met \omterm{def:g11:lenses-and-eye:defects}{ouderdomsverziendheid} heeft zijn \omterm{def:g11:lenses-and-eye:eye}{nabijpunt} op \qty{80}{cm}
maar wil op \qty{25}{cm} lezen: de bril moet de bladzijde op \qty{25}{cm}
afbeelden op een virtuele bladzijde op \qty{80}{cm}. Bereken de vereiste
\omterm{def:g11:lenses-and-eye:vergence}{lenssterkte}.
\end{exercise}

\begin{exercise}[$\star\star\star$]\label{exo:g11:lenses-and-eye:12}
Een projector moet een dia met \omterm{def:g11:lenses-and-eye:magnification}{vergroting} $\gamma = -50$ op een scherm
werpen dat \qty{5.1}{m} van de dia staat.
\begin{enumerate}
  \item Bepaal $\overline{OA}$ en $\overline{OA'}$, gegeven dat
        $\overline{OA'} - \overline{OA} = \qty{5.1}{m}$.
  \item Leid de \omterm{def:g11:lenses-and-eye:focal}{brandpuntsafstand} van het projectieobjectief af.
  \item Hoe groot is het beeld van een dia van \qty{24}{mm} hoog?
\end{enumerate}
\end{exercise}

\begin{exercise}[$\star\star\star$]\label{exo:g11:lenses-and-eye:13}
Het \omterm{def:g11:lenses-and-eye:eye}{netvlies} zit \qty{17}{mm} achter de \omterm{def:g11:lenses-and-eye:lens}{lens} van het oog.
\begin{enumerate}
  \item Bereken de \omterm{def:g11:lenses-and-eye:vergence}{lenssterkte} van het oog bij scherpstellen op oneindig
        en daarna op het \omterm{def:g11:lenses-and-eye:eye}{nabijpunt} van het \omterm{def:g11:lenses-and-eye:eye}{normale oog} (\qty{25}{cm}).
        Leid de \omterm{def:g10:signals-and-waves:amplitude}{amplitude} van de \omterm{def:g11:lenses-and-eye:eye}{accommodatie} van het \omterm{def:g11:lenses-and-eye:eye}{normale oog} af.
  \item Een ouder oog heeft zijn \omterm{def:g11:lenses-and-eye:eye}{nabijpunt} op \qty{1.0}{m}. Bereken zijn
        \omterm{def:g10:signals-and-waves:amplitude}{amplitude} en welk deel van de normale \omterm{def:g10:signals-and-waves:amplitude}{amplitude} er overblijft.
\end{enumerate}
\end{exercise}

\begin{exercise}[$\star\star\star$]\label{exo:g11:lenses-and-eye:14}
Een \omterm{def:g11:lenses-and-eye:magnifier}{loep} die als ``$\times 3$'' wordt verkocht, heeft een
\omterm{def:g11:lenses-and-eye:focal}{brandpuntsafstand} die voldoet aan $3 = 0.25/f'$ ($f'$ in \omterm{def:g10:orders-of-magnitude:unit}{meter}).
\begin{enumerate}
  \item Bereken $f'$.
  \item Waar moet een voorwerp liggen opdat zijn virtuele beeld
        \qty{25}{cm} voor de \omterm{def:g11:lenses-and-eye:lens}{lens} ligt? Bereken $\gamma$ en vergelijk met
        de aangeprezen $\times 3$.
\end{enumerate}
\end{exercise}

\begin{exercise}[$\star\star\star$]\label{exo:g11:lenses-and-eye:15}
Een \omterm{def:g11:lenses-and-eye:images}{reëel} voorwerp staat op afstand $d > 0$ voor een \omterm{def:g11:lenses-and-eye:lens}{convergerende lens}
($\overline{OA} = -d$, $f' > 0$). Toon aan dat
$\overline{OA'} = \frac{f'd}{d - f'}$ en $\gamma = \frac{f'}{f' - d}$.
Leid daaruit af: \omterm{def:g11:lenses-and-eye:images}{reëel} en omgekeerd wanneer $d > f'$; \omterm{def:g11:lenses-and-eye:images}{virtueel}, rechtop
en \emph{vergroot} ($\gamma > 1$) wanneer $0 < d < f'$ --- elke
\omterm{def:g11:lenses-and-eye:lens}{convergerende lens} is een \omterm{def:g11:lenses-and-eye:magnifier}{loep}, mits dicht genoeg gebruikt.
\end{exercise}

\section{Probleem: Eén formule, drie instrumenten}

\begin{problem}[{Weekendopgave --- dezelfde lenzenformule stelt een
fototoestel scherp, schrijft een leesbril voor, ontwerpt een loep en meet
de zoom die in je eigen oog is ingebouwd}]
\label{pb:g11:lenses-and-eye:1}
Een fotograaf draait aan een objectief van \qty{50}{mm}, een grootmoeder
houdt de krant op armlengte, een juwelier buigt zich over een diamant:
één regel algebra, $1/\overline{OA'} - 1/\overline{OA} = 1/f'$
(\cref{thm:g11:lenses-and-eye:conjugate}), bedient alle drie de
ambachten --- en daarna het oog zelf.

\medskip
\noindent\textbf{Deel I --- Eén formule.} Een \omterm{def:g11:lenses-and-eye:lens}{convergerende lens} heeft
$f' = \qty{20}{cm}$.
\begin{enumerate}
  \item Een voorwerp op \qty{60}{cm}: bereken $\overline{OA'}$ en
        $\gamma$; beschrijf het beeld.
  \item Zet het voorwerp op \qty{30}{cm} en reken opnieuw. Wat is er met
        vraag 1 verwisseld? (Lichtwegen zijn omkeerbaar.)
  \item Zet het op \qty{10}{cm}, binnen $f'$: reken opnieuw en beschrijf.
  \item Vat samen: hoe beweegt het beeld en hoe verandert zijn aard
        terwijl het voorwerp van heel ver naar $F$ schuift en daarna
        voorbij $F$?
  \item Toon aan dat een heel ver voorwerp ($1/\overline{OA} \approx 0$)
        altijd in het brandvlak wordt afgebeeld: $\overline{OA'} = f'$.
\end{enumerate}

\medskip
\noindent\textbf{Deel II --- De fotograaf.}
Een objectief heeft $f' = \qty{50}{mm}$; de sensor, \qty{24}{mm} hoog,
kan ten opzichte van de \omterm{def:g11:lenses-and-eye:lens}{lens} worden verschoven.
\begin{enumerate}[resume]
  \item Waar zit de sensor bij scherpstellen op een ver landschap?
  \item Nu een gezicht op \qty{1.0}{m}: bereken de nieuwe afstand
        lens--sensor en de verplaatsing vanaf de landschapsstand.
  \item Het mechanisme laat hoogstens \qty{55}{mm} afstand tussen \omterm{def:g11:lenses-and-eye:lens}{lens} en
        sensor toe. Bereken de kleinste scherpstelafstand.
  \item Bereken $\gamma$ bij die kortste scherpstelling. Past het beeld
        van een gezicht van \qty{20}{cm} op de sensor?
  \item ``Macro'' op ware grootte betekent $\gamma = -1$: waar moeten
        voorwerp en sensor dan staan, en waarom vraagt dat om een
        speciaal objectief?
\end{enumerate}

\medskip
\noindent\textbf{Deel III --- De grootmoeder.}
Met de jaren schuift het \omterm{def:g11:lenses-and-eye:eye}{nabijpunt} weg: dat van grootmoeder ligt op
\qty{1.0}{m}.
\begin{enumerate}[resume]
  \item Verklaar de krant op armlengte.
  \item Om op \qty{25}{cm} te lezen moet haar bril een bladzijde op
        \qty{25}{cm} omzetten in een \omterm{def:g11:lenses-and-eye:images}{virtueel beeld} op \qty{1.0}{m}:
        bereken de vereiste \omterm{def:g11:lenses-and-eye:vergence}{lenssterkte}.
  \item Bereken $\gamma$. Het beeld is vier keer groter en oogt toch niet
        groter: verklaar dat, met een vergelijking van de afmetingen
        \emph{en} de afstanden.
  \item Toon aan dat de \omterm{def:g10:pressure:pressure}{druk} alleen tussen \qty{25}{cm} en
        $f' \approx \qty{33}{cm}$ scherp is. En wat gebeurt er met een
        \emph{ver} voorwerp dat door de bril wordt bekeken?
  \item Leid in één zin af waarom er bifocale glazen bestaan.
\end{enumerate}

\medskip
\noindent\textbf{Deel IV --- De juwelier, en de zoom van het oog zelf.}
De \omterm{def:g11:lenses-and-eye:magnifier}{loep} van de juwelier heeft $f' = \qty{5.0}{cm}$; zijn \omterm{def:g11:lenses-and-eye:eye}{nabijpunt} ligt
op \qty{25}{cm}; zijn \omterm{def:g11:lenses-and-eye:eye}{netvlies} zit \qty{17}{mm} achter de \omterm{def:g11:lenses-and-eye:lens}{lens} van zijn
oog.
\begin{enumerate}[resume]
  \item Een steen \qty{4.0}{cm} onder de \omterm{def:g11:lenses-and-eye:magnifier}{loep}: plaats, aard, \omterm{def:g11:lenses-and-eye:magnification}{vergroting}?
  \item Waar moet de steen liggen opdat zijn beeld op zijn \omterm{def:g11:lenses-and-eye:eye}{nabijpunt}
        valt, en hoe groot is $\gamma$ dan?
  \item Waarom leggen ervaren juweliers de steen liever precies in $F$
        (beeld op oneindig)?
  \item Nu het oog alleen: bereken zijn \omterm{def:g11:lenses-and-eye:vergence}{lenssterkte} bij een ver voorwerp
        en daarna bij een voorwerp op \qty{25}{cm}: hoeveel \omterm{def:g11:lenses-and-eye:vergence}{dioptrie}
        \omterm{def:g11:lenses-and-eye:eye}{accommodatie} is dat?
  \item Zet beide sterkten om in \omterm{def:g11:lenses-and-eye:focal}{brandpuntsafstanden}: met hoeveel
        millimeter en hoeveel procent verandert het oog zijn eigen
        \omterm{def:g11:lenses-and-eye:focal}{brandpuntsafstand} --- de zoom die het fototoestel met
        verplaatsing en grootmoeder met $+3$~D deed?
\end{enumerate}
\end{problem}
